\subsection{Fourier 变换的定义}

定义 3.　设 \(\mu \in \mathcal{M}^{1}(\mathrm{G})\) 为 G 上的一个有界复测度。\(\mu\) 的 Fourier 变换是函数 \(\mathcal{F}_{\mathrm{G}}(\mu)\)，它定义在 \(\widehat{\mathrm{G}}\) 上，由下式给出：

\[
\mathcal {F} _ {\mathrm{G}} (\mu) (\widehat {x}) = \int_ {\mathrm{G}} \overline {{\langle \widehat {x} , x \rangle}} d \mu (x).\tag{5}
\]

我们把 \(\mu\) 的 Fourier 余变换称为函数 \(\overline{\mathcal{F}}_{\mathrm{G}}(\mu)\)，它定义在 \(\widehat{\mathrm{G}}\) 上，由下式给出：

\[
\overline {{\mathcal {F}}} _ {\mathrm{G}} (\mu) (\widehat {x}) = \int_ {\mathrm{G}} \langle \widehat {x}, x \rangle d \mu (x).\tag{6}
\]

当所考虑的群 G 没有歧义时，我们也写作 \(\mathcal{F}(\mu)\) 与 \(\overline{\mathcal{F}}(\mu)\)。我们有时也记 \(\widehat{\mu} = \mathcal{F}_{\mathrm{G}}(\mu)\)。

命题 3.　对每个测度 \(\mu \in \mathcal{M}^{1}(\mathrm{G})\)，函数 \(\mathcal{F}_{\mathrm{G}}(\mu)\) 与 \(\overline{\mathcal{F}}_{\mathrm{G}}(\mu)\) 连续且有界。映射 \(\mathcal{F}_{\mathrm{G}}: \mu \mapsto \mathcal{F}_{\mathrm{G}}(\mu)\) 与 \(\overline{\mathcal{F}}_{\mathrm{G}}: \mu \mapsto \overline{\mathcal{F}}_{\mathrm{G}}(\mu)\) 是从对合代数 \(\mathcal{M}^{1}(\mathrm{G})\) 到 \(\widehat{\mathrm{G}}\) 上有界连续函数所成的对合代数（I，第 99 页，例 1）中的连续态射。

设 \(\mu \in \mathcal{M}^1 (\mathrm{G})\)。对每个 \(\chi \in \widehat{\mathrm{G}}\) 有

\[
| \mathcal {F} (\mu) (\chi) | = \left| \int_ {\mathrm{G}} \overline {{\langle \chi , x \rangle}} d \mu (x) \right| \leqslant \| \mu \| _ {1},\tag{7}
\]

因此 \(\mu\) 的 Fourier 变换是有界的。类似地，可以验证 \(\overline{\mathcal{F}}(\mu)\) 是有界的。

若 \(\chi\) 趋于 \(\chi_0\)（于 \(\widehat{\mathbf{G}}\) 中），则函数 \(\chi\)（在 \(\mathbf{G}\) 上）在每个紧集上一致趋于 \(\chi_0\)，同时保持被常函数 1 控制，而 1 属于 \(\mathrm{L}^1 (\mathrm{G},\mu)\)。由 Lebesgue 定理（INT，IV，§3，第 7 节，定理 6），由此得 \(\mathcal{F}(\mu)(\chi)\) 趋于 \(\mathcal{F}(\mu)(\chi_0)\)。因此 \(\mathcal{F}(\mu)\) 连续。对 \(\overline{\mathcal{F}} (\mu)\) 同样成立。

对每个 \(\chi \in \widehat{\mathbf{G}}\)，映射 \(\mu \mapsto \chi(\mu) = \int \langle \chi, x \rangle d\mu(x)\) 是 \(\mathcal{M}^1(\mathbf{G})\) 的一个 Hermite 特征标（II，第 202 页，引理 1）。这蕴含 \(\mathcal{F}\) 与 \(\overline{\mathcal{F}}\) 是从 \(\mathcal{M}^1(\mathbf{G})\) 到 \(\widehat{\mathbf{G}}\) 上有界连续函数所成的对合代数中的对合代数态射。不等式 (7) 表明这些态射连续。

G 的 Fourier 变换（相应地，G 的 Fourier 余变换）是映射 \(\mu \mapsto \mathcal{F}_{\mathrm{G}}(\mu)\)（相应地，映射 \(\mu \mapsto \overline{\mathcal{F}}_{\mathrm{G}}(\mu)\)），它从 \(\mathcal{M}^1 (\mathrm{G})\) 到 \(\mathcal{C}_b(\widehat{\mathrm{G}})\) 中。

我们记下一些有用的公式，对 \(\mu\in\mathcal{M}^{1}(G)\)、\(x\in G\) 与 \(\chi\in\widehat{G}\)：

\begin{enumerate}
\def\labelenumi{(\arabic{enumi})}
\setcounter{enumi}{7}
\tightlist
\item
\end{enumerate}

\[
\overline {{\mathcal {F}}} (\mu) (\chi) = \mathcal {F} (\mu) (\chi^ {- 1}) = \overline {{\mathcal {F} (\overline {{\mu}}) (\chi)}},\tag{9}
\]

\[
\| \mathcal {F} (\mu) \| _ {\infty} = \| \overline {{\mathcal {F}}} (\mu) \| _ {\infty} \leqslant \| \mu \| _ {1},\tag{10}
\]

\[
\left\{ \begin{array}{l} \mathcal {F} (\varepsilon_ {x}) (\chi) = \overline {{\langle \chi , x \rangle}}, \\ \overline {{\mathcal {F}}} (\varepsilon_ {x}) (\chi) = \langle \chi , x \rangle , \end{array} \right.
\]

（特别地 \(\mathcal{F}(\varepsilon_e) = \overline{\mathcal{F}} (\varepsilon_e) = 1\)），

\begin{enumerate}
\def\labelenumi{(\arabic{enumi})}
\setcounter{enumi}{10}
\tightlist
\item
\end{enumerate}

\[
\left\{ \begin{array}{l} \mathcal {F} (\varepsilon_ {x} * \mu) (\chi) = \overline {{\langle \chi , x \rangle}} \mathcal {F} (\mu) (\chi), \\ \overline {{\mathcal {F}}} (\varepsilon_ {x} * \mu) (\chi) = \langle \chi , x \rangle \overline {{\mathcal {F}}} (\mu) (\chi), \end{array} \right.\tag{12}
\]

\[
\left\{ \begin{array}{l} \mathcal {F} (\chi \cdot \mu) = \varepsilon_ {\chi} * \mathcal {F} (\mu), \\ \overline {{\mathcal {F}}} (\chi \cdot \mu) = \varepsilon_ {\chi^ {- 1}} * \overline {{\mathcal {F}}} (\mu). \end{array} \right.
\]

公式 (8)、(9)、(10) 与 (11) 由定义直接得出。我们来证明公式 (12) 中的第一个，第二个是类似的。对一切 \(\xi\)（于 \(\widehat{G}\) 中），等式

\[
\begin{array}{c} \mathcal {F} (\chi \cdot \mu) (\xi) = \int_ {\mathrm{G}} \overline {{\langle \xi , x \rangle}} \langle \chi , x \rangle d \mu (x) = \int_ {\mathrm{G}} \overline {{\langle \xi \chi^ {- 1} , x \rangle}} d \mu (x) \\ = \mathcal {F} (\mu) (\xi \chi^ {- 1}) = (\varepsilon_ {\chi} * \mathcal {F} (\mu)) (\xi). \end{array}
\]

此外注意，对所有 \(\chi\in\widehat{\mathrm{G}}\) 以及所有测度 \(\mu\) 与 \(\nu\)（在 \(\mathcal{M}^{1}(\mathrm{G})\) 中），有

\[
(\varepsilon_ {\chi} * \mathcal {F} (\mu)) (\varepsilon_ {\chi} * \mathcal {F} (\nu)) = \varepsilon_ {\chi} * (\mathcal {F} (\mu) \mathcal {F} (\nu)),\tag{13}
\]

因为这个等式两边都是 \(\widehat{G}\) 上的函数，其在 \(\xi\in\widehat{G}\) 处的值为

\[
\mathcal {F} (\mu) (\xi \chi^ {- 1}) \mathcal {F} (\nu) (\xi \chi^ {- 1}).
\]

设 H 为局部紧交换群，\(\varphi\colon\mathrm{G}\to\mathrm{H}\) 为一个连续态射。设 \(\mu\in\mathcal{M}^{1}(\mathrm{G})\)。像测度 \(\varphi(\mu)\) 已有定义（INT，V，§6，第 1 节，注 1），并且由此可得 \(\mathcal{F}_{\mathrm{H}}(\varphi(\mu))=\mathcal{F}_{\mathrm{G}}(\mu)\circ\widehat{\varphi}\)（参见 INT，V，§6，第 4 节，命题 7）。

通过限制到子代数 \(L^{1}(G)\)（\(\mathcal{M}^{1}(G)\) 的），我们得到 \(L^{1}(G)\) 上 Fourier 变换与 Fourier 余变换的定义。因此，对 \(f \in L^{1}(G)\) 与 \(\chi \in \widehat{G}\) 有：

\[
\mathcal {F} _ {\mathrm{G}} (f) (\chi) = \int_ {\mathrm{G}} \overline {{\langle \chi , x \rangle}} f (x) d x, \quad \overline {{\mathcal {F}}} _ {\mathrm{G}} (f) (\chi) = \int_ {\mathrm{G}} \langle \chi , x \rangle f (x) d x.\tag{14}
\]

特别地，\(\mathcal{F}_{\mathrm{G}}(f) = \overline{\overline{\mathcal{F}}_{\mathrm{G}}(\overline{f})}\)。我们还有

\[
\overline {{\mathcal {F}}} (f) (\chi) = \chi (f)\tag{15}
\]

对所有 \(f \in \mathrm{L}^{1}(\mathrm{G})\) 与每个 \(\chi \in \widehat{G}\) 成立。

设 \(\sigma\) 为 G 的一个自同构，\(\Delta\) 为 \(\sigma\) 的模（INT，VII，§1，第 4 节，定义 4）。对 \(f \in \mathrm{L}^1(\mathrm{G})\) 有

\[
\mathcal {F} (f \circ \sigma) = \Delta^ {- 1} \mathcal {F} (f) \circ \widehat {\sigma} ^ {- 1}\tag{16}
\]

（参见同前，公式 (31)）。

若把 \(\widehat{G}\) 与 \(X(L^{1}(G))\) 等同起来（II，第 202 页，命题 1），则 Fourier 余变换不是别的，正是 Banach 代数 \(L^{1}(G)\) 的 Gelfand 变换（I，第 7 页，定义 5）。

命题 4.　Fourier 变换与 Fourier 余变换都是从 L \(^{1}\) (G) 到代数 \(\mathcal{C}_{0}(\widehat{\mathrm{G}})\)（\(\widehat{\mathrm{G}}\) 上在无穷远处为 0 的连续函数所成的代数）中的单射对合代数态射。

Fourier 余变换是从 \(L^{1}(G)\) 到 \(\widehat{G}\) 上有界连续函数代数中的一个对合代数态射。由于它与 Gelfand 变换等同，其像含于 \(\mathcal{C}_{0}(\widehat{\mathrm{G}})\)（I，第 37 页，命题 5），其核是 \(L^{1}(G)\) 的根（I，第 38 页，命题 8），而根为零（I，第 126 页，命题 22 的推论）。

我们稍后将看到（II，第 221 页，命题 13 的推论），\(\mathcal{M}^{1}(\mathrm{G})\) 上的 Fourier 余变换也是单射的。

注.　空间 \(L^{1}(G)\) 上的 Fourier 变换依赖于 Haar 测度 dx 的选取，这与 \(\mathcal{M}^{1}(G)\) 上的 Fourier 变换不同。若把 dx 换成测度 \(a \cdot dx\)（其中 \(a > 0\)），则对 G 上任何可积函数 f，f 的 Fourier 变换是 \(a\widehat{f}\)，其中 \(\widehat{f}\) 是关于测度 dx 定义的 Fourier 变换。

考虑群 G 的星代数 Stell(G)（I，第 125 页，定义 9），并把 L \(^{1}\) (G) 等同于 Stell(G) 的一个稠密子代数（I，第 126 页，命题 22）。

命题 5.　由连续性，Fourier 变换与 Fourier 余变换都唯一地扩张为从 Stell(G) 到 \(\mathcal{C}_{0}(\widehat{\mathrm{G}})\) 上的星代数同构。

Fourier 余变换由连续性扩张为从 Stell(G) 到 \(\mathcal{C}_{0}(\widehat{\mathrm{G}})\) 中的一个星代数态射。若把 \(\widehat{G}\) 与 \(\mathrm{X}(\mathrm{Stell}(\mathrm{G}))\) 等同起来（II，第 204 页，推论 1 与 II，第 202 页，命题 1），这个扩张就是 Stell(G) 的 Gelfand 变换。由 I，第 108 页，定理 1，它是一个同构。关于 Fourier 变换的断言由此得出。

我们总是用 \(\overline{\mathcal{F}}\) 与 \(\mathcal{F}\) 表示命题 5 中的同构。

推论.　\(L\)，即 \(\mathrm{L}^{1}(\mathrm{G})\) 在 G 的 Fourier 变换下的像，在 \(\mathcal{C}_{0}(\widehat{\mathrm{G}})\) 中稠密。

由于 \(L^{1}(G)\) 在 \(\text{Stell}(G)\) 中稠密，这由命题 5 得出。

命题 6.　假设 G 是紧的。规范化 Haar 测度 dx 属于 \(\mathcal{M}^{1}(\mathrm{G})\)，其 Fourier 变换是 \(\varphi_{e}\)，即 \(\{e\}\) 的特征函数。

设 \(\chi \in \widehat{\mathrm{G}}\)。由于 \(\varepsilon_y * dx = dx\) 对所有 \(y \in \mathrm{G}\) 成立，有

\[
\mathcal {F} (d x) (\chi) = \overline {{\langle \chi , y \rangle}} \mathcal {F} (d x) (\chi)
\]

（公式 (11)）。若 \(\chi \neq 1\)，则存在 \(y \in \mathrm{G}\) 使 \(\langle \chi, y \rangle \neq 1\)，从而 \(\mathcal{F}(dx)(\chi) = 0\)。若 \(\chi = 1\)，则 \(\mathcal{F}(dx)(\chi) = \int_{\mathrm{G}} dx = 1\)，因为测度 \(dx\) 是规范化的。

